\documentclass[10pt,a4paper]{amsart}
\usepackage[margin=19mm]{geometry}
\usepackage{amsmath,amssymb,mathtools}
\usepackage[scr=rsfs,cal=euler]{mathalfa}
\usepackage{xfrac}
\usepackage[unicode,hidelinks,bookmarks=false]{hyperref}
\newcommand{\ie}{i.e.\ }
\newcommand{\cf}{cf.\ }
\newcommand{\resp}{respectively\ }
\newcommand{\Subsection}{\subsection}
\providecommand{\nobreakdash}{\nobreak\hbox{-}}
\setlength{\emergencystretch}{2em}
\multlinegap0pt
\usepackage{bm}
\usepackage{bbm}
\usepackage{dsfont}
\usepackage{xspace}
\usepackage{cancel}
\usepackage{mathtools}
\usepackage{amsthm}
\makeatletter
\@ifundefined{theorem}{\newtheorem{theorem}{Theorem}}{}
\@ifundefined{lemma}{\newtheorem{lemma}[theorem]{Lemma}}{}
\makeatother
\def\FS{{\mathscr F_s}}
\def\F{{\mathscr{F}}}
\def\N{{\mathbb N}}
\title[On stronger forms of Devaney chaos]{On stronger forms of Devaney chaos}
\author{Shital H. Joshi, Ekta Shah}
\date{Source: arXiv:2505.15286v1 (CC BY 4.0)}
\begin{document}
\maketitle

\noindent\emph{Source and licence.} On stronger forms of Devaney chaos. Version arXiv:2505.15286v1 (CC BY 4.0), distributed under the Creative Commons Attribution 4.0 International licence, as recorded in the arXiv licence metadata for this exact version and shown on the abstract page https://arxiv.org/abs/2505.15286v1. Full rights notice: licenses/arxiv-2505.15286-CC-BY-4.0.txt.\par\smallskip

\noindent\textit{Editorial scope.} Counted unit: the Theorem whose source label is \texttt{T35} in the article (source0.tex lines 372--374, source label \texttt{T35}), delivered together with the article's own complete original proof of it (source0.tex lines 376--417, one \texttt{proof} environment of 42 lines). No mathematics is written, repaired or re-derived: statement and proof are the article's own text line for line. Required context retained uncounted: T6 (lines 189--191); fnf (lines 231--233). Every definition, display and label of those lines is retained unchanged. Omitted from this package and not redistributed: 1--188, 192--230, 234--371, 375--375, 418--777. Nothing inside a retained excerpt is omitted or altered: every retained body is delivered exactly as the article's lines are written, so no omission receipt of any kind accompanies this package. Citations: the retained excerpts carry no citation command at all, so no bibliography entry of the article is transcribed and no reference is redistributed. Reference adaptation: none was needed and none was made; every reference inside the delivered unit resolves inside it, so no unresolved reference or citation is produced. Printed slips kept literal: no printed word, symbol, hypothesis or number of the counted statement or of its proof was altered. Layout and class: the article's own class is \texttt{amsart} with packages including tikz, amssymb, graphicx, amsfonts, xy, amsthm, amstext, amsmath, amscd, eucal, url; the wrapper is the lane's pinned \texttt{amsart} editorial wrapper at 10pt on A4, which restates the theorem-like environments the retained text needs and adds the article's own macro definitions that the retained text needs (\texttt{\textbackslash F}, \texttt{\textbackslash FS}, \texttt{\textbackslash N}), with extra wrapper packages (bm, bbm, dsfont, xspace, cancel, mathtools). The delivered numbering is the unit's own. Assets: no retained span contains a figure environment, a graphics-inclusion command, a PDF-page inclusion, a table or a TikZ picture; no image file is copied into this package, the package declares no asset dependency and no image file is redistributed with it. Licence: version arXiv:2505.15286v1 of this article is distributed under the Creative Commons Attribution 4.0 International licence, as recorded in the arXiv licence metadata for this exact version and shown on the abstract page https://arxiv.org/abs/2505.15286v1. A full rights notice is delivered at licenses/arxiv-2505.15286-CC-BY-4.0.txt. Characters: every retained excerpt is pure ASCII, so the delivered engine, XeLaTeX with its default Latin Modern text font, reports no missing character.
\begin{lemma} \label{T6}
	Let $(X,f)$ be a dynamical system. If $f$ is uniformly continuous and $\F-$sensitive then for any $m\in \N$, $f^m$ is $\F-$sensitive, where $\F=\F_s$, $\F_t$, $\F_{ts}$ or $\F_{cf}$.
\end{lemma}

\begin{lemma}\label{fnf}
	Let $(X,f)$ be a dynamical system such that $f$ is uniformly continuous. Suppose for some $n\in \N$, $f^n$ is $\F-$sensitive. Then $f$ is $\F-$sensitive, where $\F=\F_s$, $\F_t$, $\F_{ts}$ or $\F_{cf}$.
\end{lemma}

\begin{theorem} \label{T35}
	Let  $X$ be an infinite metric space without isolated points and $f: X \longrightarrow X$ be a uniformly continuous map. Suppose $f$ is $\F-$transitive and $P(f)$ is dense in $X$.  Then $f$ is $\F-$sensitive and hence $\F-$Devaney chaotic, where $\F=\F_{s}$, $\F_t$, $\F_{ts}$ or $\F_{cf}$. 
\end{theorem}	

\begin{proof}
	In view of Lemma \ref{fnf}, it is sufficient to prove that $f^p$ is $\F-$sensitive for some $p \in \N$. Since $P(f)$ is dense in $X$ and $X$ is infinite, it follows that $P(f)$ is infinite. Let $u$ and $v$ be two periodic points such that $O_f(u)\cap O_f(v)=\phi$ and $8\delta =min\left\{d\left(f^m(u), f^n(v)\right) \; : \; m, n \geq 0 \right\}$. Then $\delta > 0$ and for each $t\in X$, $d\left(t, f^k(q)\right)\geq\delta/2$, for all $k\geq 0$. Here either $q=u$ or $q=v$.
	
	\medskip
	\noindent	Let $x \in X$ and $W$ be a neighbourhood of $x$. Put $U=W\cap B(x,\delta)$. Since $P(f)$ is dense in $X$, it follows that there is a periodic point $y$, say of period $n$, such that $y \in U$. Set $\eta = \delta/8$. Then we show that $N_{f^n}(W,\eta)=\{k\in \N: \mbox{there exists $u$, $v\in W$ with } d\left((f^n)^k(u),(f^n)^k(v)\right)>\eta\}$ contains an element of $\F$. 
	
	\smallskip
	\noindent For $x\in X$, let $q$ be a periodic point such that $d(x,f^k(q))\geq \frac{\delta}{2}$, for all $k\geq 0$. Set $$V=\bigcap_{i=0}^{n}f^{-i}\left(B(f^i(q),\eta)\right).$$ Then $V$ is a non--empty open subset of $X$ by the choice of $\eta$. 
	
	\medskip
	\noindent Since $f$ is $\mathscr{F}-$transitive, for this $U,V$, it follows that $N_f(U,V)\in \mathscr{F}$. Therefore for each $m\in N_f(U,V)$ there exists $t\in U$ such that $f^m(t)\in V$. Put 
	$$A=\left\{j\in \N:1\leq nj-m\leq n, \; m\in N_f(U,V)\right\}.$$ 
	Then $A$ is infinite as $N_f(U,V)$ is infinite. Let $j\in A$. Then $1\leq nj-m\leq n$. Further,
	$$f^{nj}(t)\in f^{nj-m}(V)\subseteq B(f^{nj-m}(q),\eta).$$
	Also, $f^{nj}(y)=y$ implies $ d(f^{nj}(t),f^{nj}(y))>2\eta $ as $y\in B(x,\eta)$ and $f^{nj}(t)\in B(f^{nj-m}(q),\eta)$. This further implies $d(f^{nj}(y),f^{nj}(x))>\eta$ or $d(f^{nj}(x),f^{nj}(t))>\eta$. Therefore $j\in N_{f^n}(W,\eta)$ and hence $A\subset N_{f^n}(W,\eta)$. We complete the proof by showing that $A\in \F$. Consider the following five cases.
	
	\medskip
	\noindent \textbf{Case 1:} $\F=\FS$
	
	\noindent Here $N_f(U,V)\in \F_s$. With the similar arguments in Case 1 of Lemma \ref{T6} one can show that $A\in \F_s$.  
	
	\medskip
	\noindent \textbf{Case 2:} $\F=\F_t$
	
	\noindent Let $k\in \N$. Since $N_f(U,V)\in \F_t$, it follows that there is a block of length $kn$ in $N_f(U,V)$. Therefore, there exists $a\in N_f(U,V)$ satisfying $$\{a,a+1,\cdots, a+n, \cdots, a+2n,\cdots , a+kn\}\subset N_f(U,V).$$
	Further, there exists $r\in \mathbb{N}$ such that $\frac{\displaystyle{a}}{\displaystyle{n}}< r \leq \frac{\displaystyle{a}}{\displaystyle{n}}+1.$ But this implies $r\in A$. Next, $\frac{\displaystyle{a+n}}{\displaystyle{n}}=\frac{\displaystyle{a}}{\displaystyle{n}}+1$ and $\frac{\displaystyle{a+n}}{\displaystyle{n}}+1=\frac{\displaystyle{a}}{\displaystyle{n}}+2$ implies
	$\frac{\displaystyle{a}}{\displaystyle{n}}+1 < r+1 \leq \frac{\displaystyle{a}}{\displaystyle{n}}+2$, which further implies $r+1\in A$. For $0\leq i<k$, $r+i\in A$. Hence $\{r,r+1,\dots,r+k-1\}\subset A$. But $k\in \N$ is arbitrary. Thus, $A\in \F_t$.
	
	\medskip
	\noindent \textbf{Case 3:} $\F=\F_{ts}$
	
	\noindent Here $N_f(U,V)\in \F_{ts}$. With the similar arguments in Case 3 of Lemma \ref{T6} one can show that $A\in \F_{ts}$. 
	
	
	\medskip
	\noindent \textbf{Case 4:} $\F=\F_{cf}$ 
	
	\noindent By Case 2, for given $a+in\in N_f(U,V)$, $r+i\in B$. Since $N_f(U,V)\in \F_{cf}$ it follows that $A\in \F_{cf}$.
	
	\smallskip
	\noindent Thus in any case $A\in \F$. Therefore $N_{f^n}(U,\eta)\in \F$. Hence, $f$ is $\F-$sensitive and therefore $\F-$Devaney chaotic.
\end{proof}

\end{document}
